\begin{textsample}{POS Dim 2 – mg – Score 23.00 – mg/mg\_em\_89.txt}  \label{ex:f2_pos_200}
4.7 \textbf{Itinerários} \textbf{formativos} - \textbf{Eletivas} Um dos componentes dos \textbf{Itinerários} \textbf{Formativos}, ao lado dos \textbf{itinerários} de aprofundamento e do Projeto de Vida, são as \textbf{Eletivas}, consideradas a parte mais \textbf{flexível} do \textbf{currículo}.
Elas foram pensadas para ampliar e diversificar os conhecimentos do estudante e ampliar seus horizontes, possibilitando maior experimentação, sem estar atreladas, necessariamente, à área de aprofundamento que ele escolher.
Ou seja, a flexibilidade possibilita a escolha de um tema associado à mesma Área do Conhecimento em que estiver se aprofundando ou optar por outro que lhe interesse associado às demais Áreas do Conhecimento.
Mesmo permitindo tratar questões diversas como música, dança, cinema, teatro, outros tipos de arte, meio ambiente, história, entre outras, de caráter mais lúdico, a \textbf{eletiva} deve estar conectada com os objetivos e os componentes do \textbf{currículo}, assim como com os eixos \textbf{estruturantes} e as Competências Gerais da BNCC.
Dessa forma, seguem as orientações gerais acima referidas na descrição dos \textbf{itinerários} de aprofundamento.
Por exemplo, uma escola de uma das inúmeras cidades históricas de Minas pode explorar em uma \textbf{eletiva} o estudo do barroco mineiro, comparando -o com as manifestações em outros países, suas origens, suas principais características.
Pode estudar a história e \textbf{trajetória} dos principais artistas barrocos, relacionar essa arte com as manifestações religiosas e a influência da igreja na formação e assentamento da população naquele e em outros \textbf{municípios} da região e o papel dos escravos na construção da arte e da riqueza do período.
Em suma, além de ampliar horizontes históricos e culturais dos estudantes, ressignifica as manifestações artísticas com as quais convive e não conhece e com a história, sociologia, antropologia etc.
Perpassa os diversos componentes das áreas do conhecimento e envolve o corpo docente na discussão da interdisciplinaridade das áreas.
Pensando na educação \textbf{profissional}, essas atividades podem, ainda, contribuir na \textbf{capacitação} dos estudantes em uma FIC como o turismo.
As \textbf{eletivas} devem ter duração de um \textbf{semestre} cada, ampliando o leque de \textbf{ofertas} e permitindo que os estudantes diversifiquem e ampliem seus conhecimentos como é o propósito desse integrante dos \textbf{Itinerários} \textbf{Formativos}.
Outro ponto importante a observar é que elas sejam elaboradas pelos professores, considerando suas capacidades e interesses, em diálogo entre o corpo docente e os discentes, identificando também os interesses e necessidades desses jovens.
Considerando que a \textbf{eletiva} pode ser integrada e abordar componentes de mais de uma área de conhecimento, o ideal é que ela seja ministrada por dois ou mais professores.
Há mais duas possibilidades de \textbf{oferta} das \textbf{eletivas}.
A primeira é quando elas se vinculam à Educação Profissional EPT e podem ser \textbf{ofertadas} como FICs ( \textbf{Curso} de \textbf{Qualificação} \textbf{Profissional} ), o que permite, inclusive, aumentar a \textbf{carga} \textbf{horária} do Itinerário \textbf{Formativo} da EPT.
Nesses casos, recomenda -se que trabalhem as habilidades requeridas pelas profissões previstas no \textbf{Catálogo} Nacional de \textbf{Cursos} \textbf{Técnicos} ( CNCT ) e na Classificação Brasileira de Ocupações ( CBO ) ( MEC, 2018 ).
Quando vinculadas à EPT, as \textbf{eletivas} são incluídas na \textbf{carga} \textbf{horária} daquele \textbf{itinerário}.
A outra possibilidade é como aproveitamento da \textbf{carga} \textbf{horária} de atividades extra escola feitas pelos estudantes, o que está previsto nas novas \textbf{Diretrizes} Curriculares Nacionais para o Ensino Médio, aprovadas pelo Conselho Nacional de Educação ( CNE ) que, no § 13, \textbf{Capítulo} II da Res. no. 3 que \textbf{atualiza} as \textbf{Diretrizes} Curriculares Nacionais para o Ensino Médio \textbf{CAPÍTULO} II, Art. 17 diz que : > § 13.
As atividades realizadas pelos estudantes, consideradas parte da \textbf{carga} \textbf{horária} do ensino médio, podem ser aulas, \textbf{cursos}, estágios, \textbf{oficinas}, trabalho supervisionado, atividades de extensão, pesquisa de campo, iniciação científica, aprendizagem \textbf{profissional}, participação em trabalhos voluntários e demais atividades com intencionalidade pedagógica orientadas pelos docentes, assim como podem ser realizadas na forma presencial – mediada ou não por tecnologia – ou a distância, inclusive mediante \textbf{regime} de parceria com instituições previamente credenciadas pelo sistema de ensino.

% matched lemmas: atualizar, capacitação, capítulo, carga, catálogo, currículo, curso, diretriz, eletivo, estruturante, flexível, formativo, horário, itinerário, município, oferta, ofertar, oficina, profissional, qualificação, regime, semestre, trajetória, técnico
\end{textsample}
